
Despite 160,000+ human judgments showing substantial consistency in identifying safety violations (Cohen's κ=0.724), standard language model embeddings reveal no geometric structure separating safe from unsafe responses (Cohen's d=0.034, statistically indistinguishable from random). The systematic three-hypothesis validation revealed this disconnect stems from pretrained encoders capturing general semantic similarity rather than safety-specific features that humans consistently detect using explicit criteria.

This negative result is scientifically valuable. It establishes clear boundaries: embedding-space clustering approaches using standard pretrained models (RoBERTa, and by extension similar general-purpose encoders) are insufficient for alignment evaluation in RLHF contexts. The failure point is localized—not dataset quality (45.6\% genuine violations), not annotation consistency (κ=0.724 agreement), but representation insufficiency for safety-specific discrimination. Alignment violations are semantically diverse (toxicity, misinformation, harmful instructions) and occupy overlapping embedding regions despite being consistently distinguishable to human annotators using policy-based criteria.

The methodological contributions remain valuable independent of the negative finding. The base-rate validation protocol provides a reusable framework for auditing preference dataset quality, addressing a gap in RLHF literature where label authenticity is assumed but not empirically verified. The three-hypothesis cascade demonstrates how systematic decomposition can pinpoint failure points in complex causal chains, applicable to alignment research beyond embedding clustering. Future work can test safety-specialized representations—reward models, safety-fine-tuned encoders—while leveraging the validated dataset quality and annotation consistency results.

\textbf{Future research directions:}

\begin{enumerate}
\item \textbf{Safety-specialized encoders:} Test whether fine-tuning on safety tasks (toxicity detection, harmful content classification) induces geometric structure absent in general-purpose pretraining.
\end{enumerate}

\begin{enumerate}
\item \textbf{Reward model embeddings:} Investigate whether RLHF reward models develop geometric structure during preference learning. Bradley-Terry models trained on HH-RLHF may learn latent "unsafe" dimensions in intermediate representations.
\end{enumerate}

\begin{enumerate}
\item \textbf{Encoder-agnostic methods:} Explore graph-based or kernel methods that do not rely on embedding-space geometry. Violations may exhibit structure in relational spaces even if not in vector embeddings.
\end{enumerate}

\begin{enumerate}
\item \textbf{Cross-dataset validation:} Extend base-rate validation to WebGPT, Summarization from Feedback, and helpfulness preferences to determine which RLHF datasets contain genuine violations versus marginal preferences.
\end{enumerate}

\begin{enumerate}
\item \textbf{Mechanistic interpretability connection:} Link preference learning dynamics to geometric structure emergence. If reward models develop clustering during training, analyze what training phases induce structure and whether it correlates with downstream policy performance.
\end{enumerate}

The question is not whether alignment structure exists—the annotation consistency proves it does—but what representations reveal it. Safety distinctions operate on semantic dimensions orthogonal to masked language modeling objectives. The path forward requires specialized representations aligned with human safety criteria rather than general-purpose similarity.
